\begin{summarybox}
\textbf{一句话记忆：固定两焦点，椭圆上的点越远离焦点连线，焦张角越大；短轴端点最大。}

\textbf{必背四式：}
\[
\boxed{r_1+r_2=2a},\qquad
\boxed{r_1r_2\cos^2\frac\theta2=b^2},
\]
\[
\boxed{\tan\frac\theta2=\frac{c|y_P|}{b^2}},\qquad
\boxed{\sin\frac\theta2\le e=\frac ca}.
\]
\textbf{取等号条件：}
$P=(0,\pm b)$，且
$\displaystyle\boxed{\theta_{\max}=2\arcsin e}$。

\textbf{最快解法：}
余弦定理化为
$r_1r_2\cos^2(\theta/2)=b^2$，
结合 $r_1r_2\le a^2$，立即得
$\sin(\theta/2)\le e$。

\textbf{四种理解路线：}
代数（余弦定理和均值不等式）；
几何（内切圆、切线长和面积）；
直观（高度 $h$ 控制观察视角）；
速解（半角恒等式 $+$ 均值不等式）。

\textbf{考场识别信号：}
若题目出现“椭圆上存在 $P$，使 $\angle F_1PF_2=\theta_0$”，
立即写
$\boxed{\sin(\theta_0/2)\le e<1}$；
若题目明确“最大焦张角为 $\theta_0$”，才写
$\boxed{e=\sin(\theta_0/2)}$。

\textbf{最重要的易错点：}
$\tan(\theta/2)=c|y_P|/b^2$，
$\cot(\theta/2)=b^2/(c|y_P|)$，二者不能互换。
\end{summarybox}
